% Draft text for the functional-verification subsection. Tables are generated: python verification/paper_runs/make_tables.py
\subsection{Functional Verification}
\label{sec:verification}

We verify every generated operator against a double-precision (float64) reference computed on the same, already quantized, inputs, so that
the reported error isolates the arithmetic of the generated design from the quantization of its inputs. Each operator variant is run through
C~simulation with $N{=}100$ independent input sets drawn uniformly (and independently per input array) from $[-r,r]$; operators whose input
range is configured as non-negative draw from $[0,r]$. We report the \emph{maximum absolute error} over all outputs and samples.

\paragraph{Formats.} We compare four configurations: (i) \texttt{ap\_fixed<16,5>} used for all storage and arithmetic; (ii) the same
\texttt{<16,5>} at the design interface, with each operator storing its operands in \texttt{<24,8>} and accumulating in \texttt{<32,10>}
with round-to-nearest (both are per-operator parameters of the generator; operands are converted at the operator boundary);
(iii) \texttt{<32,10>} throughout; and (iv) IEEE single-precision float. Fixed-point types use the Vitis default truncation and wrap unless
stated; \texttt{ap\_fixed<W,I>} has $W$ total bits of which $I$ (including the sign) are integer bits.

\paragraph{Operating range.} Fixed-point arithmetic wraps on overflow, so an operator is only meaningful over a bounded input range. For each
operator and format we search for the largest range $\pm r$ for which every sampled input stays within an error bound
\textbf{[bound: to be defined by the authors]}, confirm it with the full $N$ samples (stepping down on failure), back off by one further search step (a factor of $1.19$) because the edge of
a window is statistically soft (an overflow there needs a rare input combination), and then report the error on a fresh set of $N$
samples at that range (Table~\ref{tab:verif-range}). Because the bound is format-specific, ranges are not directly comparable across
columns; Table~\ref{tab:verif-pm1} therefore reports all four formats on the same inputs in $[-1,1]$. Windows assume uniformly distributed
inputs. The weights of matmul, MHA, sliding-window attention and convolution are held at $\pm0.1$ (a typical trained-weight scale) while
the activations are scaled, and batchnorm's parameters (scale, shift, mean, variance) are drawn from $[0.25,1]$, so that the variance has a
valid, non-degenerate floor. The range of the \texttt{tanh} activation is capped (16 for column~(ii), 24 for column~(iii)) because the Vitis
\texttt{hls::tanh} fixed-point routine crashes or returns the wrong sign for isolated large arguments (from $|x|\approx20$ in \texttt{<24,8>}
and $|x|\approx26$ in \texttt{<32,10>}). For float, verification is pass/fail against a relative--absolute tolerance, so its ``range''
is the search ceiling and its error is measured at the design's nominal input range. A dagger in Table~\ref{tab:verif-pm1} marks a format and
operator for which at least one of the $N$ samples left the format's error bound.

% Auto-generated by verification/paper_runs/make_tables.py -- do not edit by hand.
% Needs: booktabs (\toprule, \midrule, \cmidrule).
\begin{table*}[t]
\centering
\small
\setlength{\tabcolsep}{4pt}
\caption{Functional verification (C simulation, $N{=}100$ uniform input samples): the largest input range $\pm r$ over which the operator stays within the format's error bound, and the maximum absolute error over the $N$ samples at that range. Ranges are the minimum, and errors the maximum, over the operator's variants. Weights of matmul, MHA, SWA and conv are held at $\pm0.1$. Float: error at the design's nominal range; ``cap'' is the search ceiling. $^{*}$: some variant had no faithful range.}
\label{tab:verif-range}
\begin{tabular}{lrrrrrrrr}
\toprule
Operator & \multicolumn{2}{c}{\texttt{<16,5>} all} & \multicolumn{2}{c}{\texttt{<16,5>} I/O, \texttt{<24,8>} op, \texttt{<32,10>} acc} & \multicolumn{2}{c}{\texttt{<32,10>} all} & \multicolumn{2}{c}{float} \\
\cmidrule(lr){2-3}\cmidrule(lr){4-5}\cmidrule(lr){6-7}\cmidrule(lr){8-9}
 & $r$ & max abs err & $r$ & max abs err & $r$ & max abs err & $r$ & max abs err \\
\midrule
activation & 2.1 & $2.26\!\times\!10^{-2}$ & 4.2 & $6.25\!\times\!10^{-2}$ & 5.4 & $5.00\!\times\!10^{-1}$ & cap & $9.09\!\times\!10^{-7}$ \\
adaptive\_avgpool & 1.2 & $4.58\!\times\!10^{-4}$ & 19 & $4.58\!\times\!10^{-4}$ & 39 & $2.24\!\times\!10^{-7}$ & cap & $7.03\!\times\!10^{-8}$ \\
batchnorm & 5.2 & $5.69\!\times\!10^{-3}$ & 5.7 & $5.02\!\times\!10^{-4}$ & 185 & $2.18\!\times\!10^{-4}$ & cap & $5.16\!\times\!10^{-7}$ \\
conv & 2.8 & $1.58\!\times\!10^{-1}$ & 2.8 & $4.88\!\times\!10^{-4}$ & 83 & $7.71\!\times\!10^{-5}$ & cap & $5.03\!\times\!10^{-6}$ \\
dot\_product & 1.1 & $1.92\!\times\!10^{-2}$ & 1.1 & $4.82\!\times\!10^{-4}$ & 6.7 & $8.72\!\times\!10^{-6}$ & cap & $1.97\!\times\!10^{-6}$ \\
dropout & 19 & $0$ & 19 & $0$ & cap & $0$ & cap & $0$ \\
elementwise\_mult & 3.4 & $4.88\!\times\!10^{-4}$ & 3.4 & $4.88\!\times\!10^{-4}$ & 19 & $2.38\!\times\!10^{-7}$ & cap & $2.98\!\times\!10^{-8}$ \\
gemm & 0.84 & $2.16\!\times\!10^{-2}$ & 0.84 & $4.88\!\times\!10^{-4}$ & 4.8 & $9.88\!\times\!10^{-6}$ & cap & $3.34\!\times\!10^{-6}$ \\
layernorm & 0.84 & $5.16\!\times\!10^{-3}$ & 4.8 & $4.89\!\times\!10^{-4}$ & 4.8 & $1.85\!\times\!10^{-6}$ & cap & $3.44\!\times\!10^{-7}$ \\
matmul & 19 & $1.12\!\times\!10^{-2}$ & 19 & $4.88\!\times\!10^{-4}$ & cap & $5.58\!\times\!10^{-6}$ & cap & $1.71\!\times\!10^{-7}$ \\
matrix\_add & 5.7 & $0$ & 5.7 & $0$ & 181 & $0$ & cap & $5.96\!\times\!10^{-8}$ \\
maxpool & 19 & $0$ & 19 & $0$ & cap & $0$ & cap & $0$ \\
mha & 0.88 & $1.00\!\times\!10^{-1}$ & 5.5 & $4.09\!\times\!10^{-3}$ & 6.6 & $2.05\!\times\!10^{-3}$ & cap & $4.63\!\times\!10^{-7}$ \\
mmv & 0.98 & $1.92\!\times\!10^{-2}$ & 0.98 & $4.88\!\times\!10^{-4}$ & 5.7 & $9.45\!\times\!10^{-6}$ & cap & $3.13\!\times\!10^{-6}$ \\
rmsnorm & 0.84 & $5.09\!\times\!10^{-3}$ & 4.8 & $4.88\!\times\!10^{-4}$ & 4.8 & $7.82\!\times\!10^{-7}$ & cap & $2.83\!\times\!10^{-7}$ \\
swa & 3.8 & $2.34\!\times\!10^{-2}$ & 19 & $3.87\!\times\!10^{-3}$ & 21 & $3.33\!\times\!10^{-4}$ & cap & $6.41\!\times\!10^{-8}$ \\
vmm & 1.4 & $5.70\!\times\!10^{-3}$ & 1.4 & $4.88\!\times\!10^{-4}$ & 7.8 & $2.77\!\times\!10^{-6}$ & cap & $7.97\!\times\!10^{-7}$ \\
\bottomrule
\end{tabular}
\end{table*}

\begin{table}[t]
\centering
\small
\caption{Maximum absolute error over $N{=}100$ samples with every input uniform in $[-1,1]$ (inputs whose configured range is non-negative: $[0,1]$), for all four formats. $^{\dagger}$: at least one sample left the format's error bound (overflow or wrap) for some variant.}
\label{tab:verif-pm1}
\begin{tabular}{lrrrr}
\toprule
Operator & \texttt{<16,5>} all & \texttt{<16,5>} I/O, \texttt{<24,8>} op, \texttt{<32,10>} acc & \texttt{<32,10>} all & float \\
\midrule
activation & $1.32\!\times\!10^{-1}$$^{\dagger}$ & $8.89\!\times\!10^{-4}$ & $6.47\!\times\!10^{-7}$ & $1.52\!\times\!10^{-7}$ \\
adaptive\_avgpool & $4.58\!\times\!10^{-4}$ & $4.58\!\times\!10^{-4}$ & $2.24\!\times\!10^{-7}$ & $7.03\!\times\!10^{-8}$ \\
batchnorm & $2.04\!\times\!10^{-3}$ & $5.02\!\times\!10^{-4}$ & $1.85\!\times\!10^{-6}$ & $5.16\!\times\!10^{-7}$ \\
conv & $1.56\!\times\!10^{-1}$ & $4.88\!\times\!10^{-4}$ & $7.79\!\times\!10^{-5}$ & $5.03\!\times\!10^{-6}$ \\
dot\_product & $1.84\!\times\!10^{-2}$ & $4.85\!\times\!10^{-4}$ & $9.01\!\times\!10^{-6}$ & $1.97\!\times\!10^{-6}$ \\
dropout & $0$ & $0$ & $0$ & $0$ \\
elementwise\_mult & $4.88\!\times\!10^{-4}$ & $4.88\!\times\!10^{-4}$ & $2.38\!\times\!10^{-7}$ & $2.98\!\times\!10^{-8}$ \\
gemm & $2.02\!\times\!10^{-2}$ & $4.88\!\times\!10^{-4}$ & $1.01\!\times\!10^{-5}$ & $3.34\!\times\!10^{-6}$ \\
layernorm & $4.70\!\times\!10^{-3}$$^{\dagger}$ & $4.88\!\times\!10^{-4}$ & $1.74\!\times\!10^{-6}$ & $3.44\!\times\!10^{-7}$ \\
matmul & $1.10\!\times\!10^{-2}$ & $4.88\!\times\!10^{-4}$ & $5.50\!\times\!10^{-6}$ & $1.71\!\times\!10^{-7}$ \\
matrix\_add & $0$ & $0$ & $0$ & $5.96\!\times\!10^{-8}$ \\
maxpool & $0$ & $0$ & $0$ & $0$ \\
mha & $1.13\!\times\!10^{-1}$ & $5.44\!\times\!10^{-4}$ & $5.62\!\times\!10^{-5}$ & $4.63\!\times\!10^{-7}$ \\
mmv & $1.93\!\times\!10^{-2}$ & $4.88\!\times\!10^{-4}$ & $9.40\!\times\!10^{-6}$ & $3.13\!\times\!10^{-6}$ \\
rmsnorm & $4.03\!\times\!10^{-3}$$^{\dagger}$ & $4.88\!\times\!10^{-4}$ & $1.77\!\times\!10^{-6}$ & $2.83\!\times\!10^{-7}$ \\
swa & $1.16\!\times\!10^{-2}$ & $5.08\!\times\!10^{-4}$ & $5.66\!\times\!10^{-6}$ & $6.41\!\times\!10^{-8}$ \\
vmm & $5.59\!\times\!10^{-3}$ & $4.88\!\times\!10^{-4}$ & $2.76\!\times\!10^{-6}$ & $7.97\!\times\!10^{-7}$ \\
\bottomrule
\end{tabular}
\end{table}

